(1) 同一内部模式时输入态为 \(a^\dagger b^\dagger|0\rangle\)。分束器使其变为 \(\tfrac12(c^\dagger+d^\dagger)(c^\dagger-d^\dagger)|0\rangle=\tfrac12[(c^\dagger)^2-(d^\dagger)^2]|0\rangle\)，因为不同输出模式的玻色产生算符彼此对易，交叉项完全相消。归一化后输出态为 \((|2_c,0_d\rangle-|0_c,2_d\rangle)/\sqrt2\)。

(2) 态中没有 \(|1_c,1_d\rangle\) 分量，因此完全同模的符合概率 \(P_{\rm coin}=0\)，而两个光子同时从 \(c\) 或 \(d\) 输出的概率各为 \(1/2\)。这不是两个单光子独立抛硬币的结果，而是两条不可区分量子路径的概率振幅先相加、再取模平方导致的玻色增强和反聚束输出。

(3) 当内部态可区分时，输出符合事件有两个物理排列，它们的干涉交叉项正比 \(|\langle f|g\rangle|^2=|\mu|^2\)。严格对未被探测的内部自由度求和得到 \(P_{\rm coin}=\frac12(1-|\mu|^2)\)。完全重叠 \(|\mu|=1\) 时恢复零符合；若两模式正交 \(\mu=0\)，干涉项消失而 \(P_{\rm coin}=1/2\)，等于两个可区分粒子独立经过平衡分束器的符合概率。

(4) 概率守恒和两端对称性给 \(P_{20}=P_{02}=(1-P_{\rm coin})/2=(1+|\mu|^2)/4\)。\(\mu=0\) 时输出分布为 \((1/4,1/2,1/4)\)，\(|\mu|=1\) 时为 \((1/2,0,1/2)\)。实验中扫描相对延迟会改变波包重叠 \(\mu(\tau)\)，符合曲线出现 HOM dip，其可见度在理想纯单光子模型中与可区分度直接相关；损耗、多光子源和暗计数会引入额外修正。
